\documentclass[12pt]{article}
\usepackage[utf8]{inputenc}
\usepackage[T1]{fontenc}
\usepackage{hyperref}
\title{Pyam Language Resources for K-12 Curriculum Development}
\author{Compiled by Claude Sonnet 5.5 (claude-sonnet-5-5) for Andr\'as Kornai}
\date{}
\begin{document}
\maketitle

\section{Introduction}
Pyam (also Fyam, Pyem; ISO 639-3 pym; Glottocode fyam1238) is a Plateau language of Plateau State, Nigeria, spoken mainly in Mangu Local Government Area around Gindiri. Speaker estimates in the sources vary widely, and shift to Hausa is reported. Search results mention that in 2018 the Ba-Pyam Azonci Association commissioned books in Pyem; the books themselves were not located.

\section{Primary Educational Resources}
No primer or textbook was located directly. The 2018 book project is mentioned on the Wikipedia page below.

\section{Linguistic and Reference Materials}
\begin{itemize}
\item Webonary dictionary site for Pyam. \url{https://www.webonary.org/pyam/overview/introduction/}
\item Max Planck Institute numeral database entry for Pyam. \url{https://lingweb.eva.mpg.de/channumerals/Pyam.htm}
\item Wikipedia: East Plateau languages (classification context). \url{https://en.wikipedia.org/wiki/East_Plateau_languages}
\end{itemize}

\section{Religious and Community Resources}
\begin{itemize}
\item Wikipedia: Pyem people (community, location, language shift, 2018 book project). \url{https://en.wikipedia.org/wiki/Pyem_people}
\end{itemize}

\section{Digital and Audio Resources}
\begin{itemize}
\item Endangered Languages Project entry for Pyam. \url{https://www.endangeredlanguages.com/lang/5303}
\end{itemize}

\section{Curriculum Development Context}
A Webonary dictionary exists, and a community association has commissioned books, which suggests local interest. No school programme was found. Contacting the community association (not verified here) would be the natural next step.

\end{document}